% Generated by scripts/paper/make_cost_table.py from measure_embedding_cost.py output.
\begin{table}[htbp]
  \centering
  \small
  \setlength{\tabcolsep}{3pt}
  \begin{tabular}{@{} l r r r r r r r @{}}
    \toprule
    & \textbf{Params} & \textbf{Weights} & & \multicolumn{3}{c}{\textbf{Molecules per second}} & \textbf{Peak RAM} \\
    \cmidrule(lr){5-7}
    \textbf{Representation} & \textbf{(M)} & \textbf{(MB)} & \textbf{Dim.} & \textbf{CPU} & \textbf{MPS 32} & \textbf{MPS 128} & \textbf{(GB)} \\
    \midrule
    ECFP4, binary & --- & --- & 2,048 & 9,944 & --- & --- & 0.37 \\
    ECFP4, count & --- & --- & 2,048 & 9,034 & --- & --- & 0.49 \\
    \model{}-small & 33.9 & 136.6 & 512 & 280 & 578 & 577 & 1.11 \\
    \model{}-base & 113.8 & 457.4 & 768 & 124 & 326 & 251 & 1.68 \\
    \bottomrule
  \end{tabular}
  \caption{Feature-extraction cost on one workstation (Apple M1 Max, 10 CPU cores, 32\,GB; PyTorch 2.12.1) for 10,000 distinct benchmark test molecules. Throughput is the median of 3 interleaved repeats after one untimed warm-up pass, and counts every input molecule. It includes SMILES parsing for ECFP4 and, for \model{}, SMILES-to-\selfies{} conversion, tokenisation, the encoder forward pass and mean pooling, with inputs capped at 128 tokens; \model{} embedded 9,696 of the molecules and its featuriser rejected the rest as lossy or unsupported inputs. CPU runs used batch size 32 and 8 PyTorch threads; the MPS columns give batch sizes 32 and 128 on the Apple GPU. \emph{Params}: encoder parameters used for embedding, excluding the MLM head; \emph{Weights}: size of the model weight file; \emph{Peak RAM}: peak resident memory of a fresh CPU process. ECFP4 uses RDKit Morgan fingerprints (radius 2, 2,048 bits) in one process. Other applications were running (one-minute load average 5--16); absolute rates depend on the machine and its load.}
  \label{tab:extraction-cost}
\end{table}
